% Ad Atticum 14.20 (ad-atticum-14-20) --- English translation
% Translated by Claude (Anthropic) from the Latin (Perseus).
% Style guide: STYLE.md.

\ciceroLetterOpener{CICERO ATTICO salutem}

\ciceroSection{1} From the Pompeian villa I was carried by ship to the hospitality of our friend Lucullus on the sixth before the Ides, at about the third hour. As I stepped off the boat I was handed your letter, which your courier was said to have brought to the Cumean place, dated the Nones of May. From Lucullus's I came on the next day, at about the same hour, to the Puteolan villa. There I received two more letters from you, one dated the Nones, the other the seventh before the Ides, from Lanuvium.

\ciceroSection{2} Listen, then, to my answers to them all. First: what you write about my affairs --- both the payment and the Albian business --- is welcome. As for your Buthrotum: while I was at the Pompeian villa, Antony came to Misenum. From there he was gone before I had even heard he was come, and into Samnium, of which you can guess what you can hope. So Buthrotum will have to be done at Rome. Lucius Antonius's harangue was a horror; Dolabella's brilliant. He may keep the money for all I care, so long as he pays out on the Ides. I am sorry Tertulla has miscarried. There need now to be Cassii to plant as much as Bruti. About the queen --- and that ``Caesar'' of hers --- I should be glad of news. I have settled the first letter; I come to the second.

\ciceroSection{3} About the Quintuses, and about Buthrotum: when I am there, as you write. Your supporting Cicero is welcome. As to your thinking me mistaken in supposing the commonwealth hangs upon Brutus --- it is so: either there will be none, or it will be saved by him, or by such as him. Now, as to your urging me to send you the speech in writing: accept from me, my dear Atticus, a general theorem (\textit{[Greek: katholikon theorema]}) about those matters in which we are sufficiently practised. No poet, no orator, has ever lived who thought any other man better than himself. This befalls even the bad ones. What then must you suppose of Brutus, with his gift and his learning? We have had recent experience of it, even, over the edict. I had drafted one at your request. Mine pleased me; his pleased him. What is more, when, almost won over by his own entreaties, I had written to him about the best style of oratory, he wrote back not only to me but to you as well that what pleased me did not please him. So leave each man, please, to write for himself; \textit{to each his own betrothed, to me my own; to each his own beloved, to me mine.} Not elegant, that. It is Atilius, the harshest of poets, who says it. And, by the gods, may it be open to that fellow to harangue the people! For if he is permitted to be in the city in safety, we have won. The leader of a new civil war, either no one will follow, or those will follow who can be easily defeated.

\ciceroSection{4} I come to the third. I am glad my letter was welcome to Brutus and Cassius. I have written back to them therefore. That they wish Hirtius to be improved through me: I am taking pains, and the man himself talks excellently, but he lives and keeps house with Balbus, who likewise talks well. What you are to believe is for you to see. I see that Dolabella greatly pleases you; me too, signally. With Pansa I spent some time at the Pompeian villa. He plainly satisfied me that he was right-minded and longed for peace. That a pretext for arms is being looked for, I can plainly see. The edict of Brutus and Cassius I approve. As for your wanting me to take up the question of what these men ought to do: counsels are matters of the hour, and you see them change from hour to hour. Dolabella, both that first action of his and now this harangue against Antony, have advanced things very considerably, it seems to me. The thing was simply moving forward. Now, however, we seem likely to have a leader; which is the one thing the country towns and the loyal men want.

\ciceroSection{5} You mention Epicurus and dare to say ``not'' (\textit{[Greek: m\=e]}). Doesn't even the slight scowl of our Brutus put you off such talk? Quintus my nephew, as you write, is Antony's little right hand. Through him, then, we shall easily carry what we want. I am waiting --- if, as you suppose, Lucius Antonius did bring Octavian forward --- to hear what kind of harangue it was. I have dashed this off; for Cassius's courier is just leaving. I was on the point of going to greet Pilia, then to a feast at Vestorius's by skiff. Very warmest greetings to Attica.
